\scanpage{51}
% loc: heading above the rule
Remark on the characterization of \deleted{\unsure{qu}}\inserted{essentially} \uline{algebraic} de Rham cohomology classes. by bigrading or by reality ....
\par\noindent\rule{\linewidth}{0.3pt}\par
% loc: opening and Hom insertion
$K\subset\mathbb C$ subfield, $M$ motive over $K$, whence $T_{\mathrm{DR}}(M)$. Let us consider $T_{\mathrm{DR}}(M\unsure{\natural|\leftrightharpoons})$\inserted{$=\operatorname{Hom}(1,M)\otimes_{\mathbb Q}K$}, one sees that it is a submodule of $T_{\mathrm{DR}}(M)$, \deleted{which is}\inserted{and that}
% loc: clauses a and b
\par a) contained in $T_{\mathrm{DR}}(M)_{\mathbb C}^{00}$.
\par b) \inserted{$\operatorname{Hom}(1,M)\otimes_{\mathbb Q}(K\cap\mathbb R)$ is} contained in $\bigl(T_{\mathrm{DR}}(M)_{\mathbb C}\bigr)_{\mathbb R}$ \deleted{[and where $\operatorname{Hom}(1,M)\otimes_{\mathbb Q}(K\cap\mathbb R)$ is contained in $\bigl(T_{\mathrm{DR}}(M)_{\mathbb C}\bigr)_{\mathbb R}$]}

% loc: question
Do we have a converse ?

% loc: item 1, opening and group insertions
$1^\circ$) Not for \deleted{a) alone, even if}\inserted{a) and b) simultaneously} $K=\mathbb Q$, if the motivic Galois group $G$ \inserted{for the category of motives $\unsure{\mathcal M|M}$ in which we work} is not connected. Indeed, if $G$ is \deleted{\unsure{discrete}} finite [Case of effective motives of weight 0] condition a) is vacuous, since $T_{\mathrm{DR}}(M)_{\mathbb C}^{00}=T_{\mathrm{DR}}(M)_{\mathbb C}$, \inserted{thus converse of a) true if and only if $G=(e)$.} on the other hand \deleted{b) is} the converse to b) is true if and only if the Galois extension $L$ with group $G$ of $\mathbb Q$ defined by $\mathcal M$ is \uline{totally real} [one tests with the $G$-module $M=\mathbb Q[C]$, which gives $T_{\mathrm{DR}}(M)=L$ \ldots] ; thus in this case one has a converse to a) and b).

% loc: connected group and embeddings
One must therefore restrict oneself to a connected $G$, which suggests looking at the converses for the case $K$ algebraically closed, \inserted{for example $K=\overline{\mathbb Q}$} One can moreover strengthen the condition, by looking at \uline{all} the embeddings $K\hookrightarrow\mathbb C$.

% loc: item 2, introductory lines
$2^\circ$) The converse to a) is false, \deleted{\unsure{even if}} (even for $\overline{\mathbb Q}$). it suffices to take a motivic Galois group $G$ which is a \uline{torus} of dimension $\geq 3$. One sees that one can find \inserted{on}
% loc: pair of arrows, target, and marginal additions
$j_1,j_2:{\mathbb G_{\mathrm m\,K}}^{\unsure{p|\rho}}\rightrightarrows G'_{\deleted{\unsure{K}}\unsure{\mathbb C|K}}$\inserted{$(=\operatorname{ad}(P)$, since $G$ \unsure{commutative!|commutative})} which define \inserted{(for every $K\longrightarrow\mathbb C$)} the bigrading of
% loc: struck margin beside the two arrows
\margin{\deleted{\uline{NB} $j_1$ and $j_2$ \unsure{come} from \unsure{a choice} of \unsure{two} \unsure{embeddings} \unsure{into} $\mathbb C$, and \illegible{1}}}
% loc: scalar extensions and torsor
$\deleted{G}V_{\mathbb C}\simeq V'\otimes_K\mathbb C$ \quad [$V'=P\mathbin{\overset{G_K}{\deleted{\unsure{\times}}}}V$, $P\unsure{\simeq|=}\underline{\operatorname{Isom}}_{\otimes}(T_{\mathrm B}\otimes K,T_{\mathrm{DR}})$]
% loc: degree-zero representation
Thus a \inserted{faithful} representation of degree 0 of $G$ gives, $\deleted{V^{G'}}$ if $T$ is the image of $(j_1,j_2):\mathbb G_{\mathrm m\,K}\times\mathbb G_{\mathrm m\,K}\mathbin{-}G'$, a $V'$ for which
% loc: final equality and inequality
${V'}^{T}\inserted{(=T_{\mathrm{DR}}\cap(T_{\mathrm{DR}}\deleted{\unsure{)}}{}_{\mathbb C}^{00})}$
\inserted{\deleted{\unsure{some}}} is $\ne {V'}^{G'}$.
% Watermark lettering supplied by C is not recoverable on this image; retained as an audit question in the changes log.
% archive sheet number: 50
